% =============================================================================
%  AMIS Addons - author guide
%  Public documentation for the AMIS addon catalogue.
%  Build: pdflatex main && pdflatex main
% =============================================================================
\documentclass[11pt,a4paper]{article}

\usepackage[T1]{fontenc}
\usepackage[utf8]{inputenc}
\usepackage{lmodern}
\usepackage[margin=2.4cm]{geometry}
\usepackage{microtype}
\usepackage{xcolor}
\usepackage{booktabs}
\usepackage{longtable}
\usepackage{listings}
\usepackage{enumitem}
\usepackage{titlesec}
\usepackage{fancyhdr}
\usepackage{caption}
\usepackage{amsmath,amssymb}
\usepackage{fancyvrb}
\usepackage[colorlinks=true,linkcolor=teal,urlcolor=teal]{hyperref}

\definecolor{teal}{HTML}{265C74}
\definecolor{amber}{HTML}{D94F2B}
\definecolor{codebg}{HTML}{F4F6F7}

\captionsetup{font=small,labelfont={bf,color=teal}}
\titleformat{\section}{\Large\bfseries\color{teal}}{\thesection}{0.7em}{}
\titleformat{\subsection}{\large\bfseries}{\thesubsection}{0.6em}{}

\pagestyle{fancy}
\fancyhf{}
\fancyhead[L]{\small\color{teal}AMIS Addons --- author guide}
\fancyhead[R]{\small\color{teal}\thepage}
\renewcommand{\headrulewidth}{0.4pt}

\lstdefinestyle{amis}{
  backgroundcolor=\color{codebg},
  basicstyle=\ttfamily\footnotesize,
  keywordstyle=\color{teal}\bfseries,
  commentstyle=\color{gray}\itshape,
  stringstyle=\color{amber},
  frame=single,
  rulecolor=\color{gray!40},
  breaklines=true,
  showstringspaces=false,
  columns=fixed,
  basewidth=0.5em,
  keepspaces=true,
}
\lstset{style=amis}

\newcommand{\file}[1]{\texttt{\detokenize{#1}}}
\newcommand{\code}[1]{\texttt{\detokenize{#1}}}
\newcommand{\ui}[1]{\textsf{#1}}

\begin{document}

\begin{titlepage}
  \centering
  \vspace*{3.5cm}
  {\Huge\bfseries\color{teal} AMIS Addons\par}
  \vspace{0.5cm}
  {\Large Author guide\par}
  \vspace{0.3cm}
  {\large How to write, package and publish a segmentation method\par}
  \vspace{1.5cm}
  {\large Addon API v1 \quad\textbullet\quad AMIS 0.1.0\par}
  \vfill
  {\small Catalogue: \url{https://github.com/ADOBApps/amis-addons}\par}
  {\small Application: \url{https://github.com/ADOBApps/advancedmedicalis}\par}
\end{titlepage}

\tableofcontents
\newpage

\section{What an addon is}

An addon is a folder that the AMIS application discovers, imports and runs. It
owns a segmentation method end to end: its model, its training recipe and its
window.

\begin{lstlisting}
my_method/
    manifest.json     identity and compatibility   (required)
    __init__.py       behaviour, exposes ADDON     (required)
    panel.json        the interface, as data       (recommended)
    model.py          torch model + build_model    (optional)
    README.md         what it does                 (optional)
\end{lstlisting}

Three rules explain the rest:

\begin{enumerate}[leftmargin=1.4em,itemsep=2pt]
  \item \textbf{The folder is the unit.} Install = copy, remove = delete.
  \item \textbf{The folder name is the id.} \file{my_method/} declares
        \code{"id": "my_method"} and is imported as \code{amis_addons.my_method}.
  \item \textbf{The window is data.} \file{panel.json} is interpreted at run time,
        so the visual designer and the running application show the same thing.
\end{enumerate}

\section{Quick start}

\subsection{In the application}

Toolbar \(\to\) \ui{New addon\dots}: choose an id, a display name and a kind
(\emph{model} or \emph{declarative}). AMIS writes a \emph{runnable} skeleton into
its workspace, opens the files in the \emph{Code} tab, and the addon already
works: \ui{Run} prints the model's parameter count.

Press \ui{Validate} (compile + JSON + manifest), then \ui{Test load} (the real
registry imports it and builds the panel), then \ui{Install}, then
\ui{Reload addons}.

\subsection{From the terminal}

\begin{lstlisting}[language=bash]
python -m amis.market validate ../my_method --load   # compile + import + panel
python -m amis.market package  ../my_method --catalogue-entry
python -m amis.market verify   my_method-0.1.0.zip
\end{lstlisting}

\subsection{The smallest complete addon}

\begin{lstlisting}[language=Python]
# my_method/__init__.py
from amis.core.addon_api import AddonBase

class MyMethod(AddonBase):
    def on_run_clicked(self) -> None:
        self.ui.set("log", "hello from my method")

ADDON = MyMethod
\end{lstlisting}

\noindent with a \file{panel.json} containing a button named \code{run}, a text
control named \code{log}, and \code{'api': 1} in the manifest. That is the whole
contract; everything below is detail.

\section{The manifest}

\begin{lstlisting}[language=Python]
{
  "id": "my_method",
  "name": "My Method",
  "version": "0.1.0",
  "api": 1,
  "author": "Your Name",
  "license": "GPL-3.0-or-later",
  "description": "One sentence.",
  "category": "segmentation",
  "main": "__init__.py",
  "capabilities": ["train", "predict", "export"],
  "order": 500,
  "tags": ["2d", "unet"],
  "config": {}
}
\end{lstlisting}

\begin{longtable}{@{}p{3.2cm}p{1.5cm}p{9.3cm}@{}}
\toprule
\textbf{Key} & \textbf{Req.} & \textbf{Meaning} \\
\midrule
\endhead
\code{id} & yes & identifier; must equal the directory name \\
\code{name} & yes & label in the method list \\
\code{api} & yes & must be \code{1} for this release \\
\code{version} & no & semantic version, e.g. \code{1.2.0} \\
\code{license} & yes* & required for publication; SPDX identifier preferred \\
\code{capabilities} & no & any of \code{train}, \code{predict}, \code{export} \\
\code{order} & no & ascending sort key (bundled: 10, 20, 900) \\
\code{config} & no & free-form; read as \code{context.manifest.config} \\
\bottomrule
\end{longtable}

\noindent *) enforced by the catalogue validator and the AI gate, not by the
loader, so an unpublished addon still runs without it.

\section{Behaviour: \file{\_\_init\_\_.py}}

\subsection{Lifecycle}

\begin{verbatim}
initialize()      once, before the panel exists; validate files, load nothing heavy
build_panel()     first time the method is selected; default: render panel.json
on_activated()    every time the tab becomes current
shutdown()        window close or reload; must be idempotent
\end{verbatim}

\code{AddonBase} implements every hook except \code{build_panel()}, which renders
\file{panel.json} when present.

\subsection{The context}

\begin{longtable}{@{}p{3.6cm}p{10.8cm}@{}}
\toprule
\textbf{Member} & \textbf{Use} \\
\midrule
\endhead
\code{context.data_dir} & writable folder for this addon (weights, exports, caches) \\
\code{context.addon_dir} & the (read-only) addon folder itself \\
\code{context.settings} & \code{get}/\code{set}/\code{update}, scoped to your addon \\
\code{context.logger} & configured logger; also written to the application log \\
\code{context.qt_parent} & parent for dialogs you open \\
\code{context.notify(msg, level)} & show a message in the host \\
\bottomrule
\end{longtable}

\subsection{Reading and writing the window}

\code{self.ui} is the renderer:

\begin{lstlisting}[language=Python]
epochs = self.ui.get("epochs")          # int
self.ui.set("progress", 40)             # int
preview = self.ui.values["preview"]     # raw widget: ImageView
preview.set_image(mask)
\end{lstlisting}

\begin{quote}
\small Handlers may run before the panel exists (tests, the CLI). Guard with
\code{if self.ui is not None}. A handler that raises inside a Qt slot can take
the application down.
\end{quote}

\subsection{Handlers}

A button calls, in this order: the panel's \code{on_click}, then
\code{on_<name>_clicked}. An input calls \code{on_change}, then
\code{on_<name>_changed} and receives the current value.

\begin{lstlisting}[language=Python]
def on_run_clicked(self) -> None: ...
def on_threshold_changed(self, value: float) -> None: ...
\end{lstlisting}

A missing handler is not an error (the control simply does nothing), which is why
the validator reports it.

\subsection{Threads}

Training and downloads run in a worker thread. Widgets must only be touched from
the GUI thread, so report through signals:

\begin{lstlisting}[language=Python]
class _Bridge(QObject):
    log = Signal(str)

def on_start_clicked(self) -> None:
    self._bridge = _Bridge()
    self._bridge.log.connect(lambda m: self.ui.set("log", m))
    self._trainer = Trainer(..., on_log=self._bridge.log.emit)
    self._trainer.start()
\end{lstlisting}

The trainer from \code{amis.ml} already follows this design: it takes plain
callables and knows nothing about Qt.

\section{The window: \file{panel.json}}

\begin{lstlisting}[language=Python]
{"amis_panel": 1, "title": "My Method", "width": 600, "height": 700,
 "root": {"type": "vbox", "children": [
   {"type": "group", "props": {"title": "Settings"}, "children": [
     {"type": "dspin", "name": "threshold",
      "props": {"label": "Threshold", "value": 0.5, "minimum": 0.0,
                "maximum": 1.0, "step": 0.05, "decimals": 2}}]},
   {"type": "button", "name": "run", "props": {"label": "Run"},
    "on_click": "on_run_clicked"},
   {"type": "image", "name": "preview", "props": {"height": 240}},
   {"type": "text", "name": "log", "props": {"text": "", "rows": 6}}]}}
\end{lstlisting}

\begin{longtable}{@{}p{2.6cm}p{3.0cm}p{8.4cm}@{}}
\toprule
\textbf{Type} & \textbf{Qt widget} & \textbf{Properties} \\
\midrule
\endhead
\code{vbox}, \code{hbox} & layout & -- (container) \\
\code{group} & \code{QGroupBox} & \code{title} (container) \\
\code{scroll} & \code{QScrollArea} & -- (container) \\
\code{form} & \code{QFormLayout} & -- (container) \\
\code{tabs} & \code{QTabWidget} & each child needs \code{props.title} \\
\code{button\_row} & horizontal group & -- (container) \\
\code{spacer} & filler & \code{height} \\
\code{separator} & line & \code{orientation} (\code{h}/\code{v}) \\
\code{label} & \code{QLabel} & \code{text}, \code{bold}, \code{wrap} \\
\code{text} & \code{QTextEdit} & \code{text}, \code{rows}, \code{readonly}, \code{mono} \\
\code{image} & image view & \code{height} \\
\code{progress} & \code{QProgressBar} & \code{minimum}, \code{maximum}, \code{value} \\
\code{line} & \code{QLineEdit} & \code{text}, \code{placeholder}, \code{label} \\
\code{spin} & \code{QSpinBox} & \code{label}, \code{value}, \code{minimum}, \code{maximum}, \code{step} \\
\code{dspin} & \code{QDoubleSpinBox} & the above plus \code{decimals} \\
\code{slider} & \code{QSlider} & \code{orientation}, range \\
\code{check} & \code{QCheckBox} & \code{text}, \code{checked} \\
\code{combo} & \code{QComboBox} & \code{items}, \code{current}, \code{label} \\
\code{file} & chooser & \code{label}, \code{mode} (\code{open}/\code{save}), \code{filter} \\
\code{dir} & folder chooser & \code{label}, \code{path} \\
\code{button} & \code{QPushButton} & \code{label}, \code{enabled} \\
\bottomrule
\end{longtable}

Rules: containers carry \code{children}, leaves never do; names are unique;
numbers always display a decimal point; there is no styling property, because
the panel follows the desktop theme.

\section{The model: \file{model.py}}

\begin{lstlisting}[language=Python]
import torch, torch.nn as nn
DEFAULT_CONFIG = {"in_channels": 1, "out_channels": 1, "base_channels": 8}

class MyNet(nn.Module):
    def __init__(self, in_channels=1, out_channels=1, base_channels=8):
        super().__init__()
        self.body = nn.Sequential(
            nn.Conv2d(in_channels, base_channels, 3, padding=1), nn.ReLU(True),
            nn.Conv2d(base_channels, out_channels, 1))
    def forward(self, x):
        return self.body(x)

def build_model(config=None):
    merged = dict(DEFAULT_CONFIG); merged.update(config or {})
    return MyNet(**merged)
\end{lstlisting}

Contract: \code{build_model(config) -> nn.Module}; output \emph{logits} shaped
\((N, C, H, W)\) with the input's spatial size; one channel for binary
segmentation. Inputs arrive padded to a multiple of 16 and are cropped back, but
a model that can handle odd sizes itself is more robust.

\section{Shared plumbing}

\begin{lstlisting}[language=Python]
from amis.ml import data as data_io          # discover_cases, load_gray, load_mask, save_mask
from amis.ml import metrics as metric_lib    # dice, iou, precision, sensitivity, all_metrics
from amis.ml.losses import bce_dice_loss
from amis.ml.trainer import Trainer, TrainingConfig
from amis.ml.backend import make_segmenter
from amis.ml.onnx_export import export_onnx
\end{lstlisting}

\begin{longtable}{@{}p{4.6cm}p{9.8cm}@{}}
\toprule
\textbf{Helper} & \textbf{What it does} \\
\midrule
\endhead
\code{data_io.discover_cases(images, masks)} & pairs images and masks by file stem \\
\code{data_io.load_gray(path)} & float32 \((H,W)\) in \([0,1]\) \\
\code{Trainer(...).start()} & threaded training; writes \code{best.pt}, \code{last.pt}, \code{history.json} \\
\code{make_segmenter(...)} & torch or ONNX predictor with padding and cropping \\
\code{metric_lib.all_metrics(p, y)} & Dice, IoU, precision, sensitivity, specificity, accuracy \\
\code{export_onnx(...)} & ONNX export plus an ONNX Runtime parity check \\
\bottomrule
\end{longtable}

Metric definitions, with \(TP,FP,FN,TN\) pixel counts and \(\varepsilon=10^{-7}\):

\begin{equation*}
\mathrm{Dice} = \frac{2TP+\varepsilon}{2TP+FP+FN+\varepsilon}, \quad
\mathrm{IoU} = \frac{TP+\varepsilon}{TP+FP+FN+\varepsilon}, \quad
\mathrm{Precision} = \frac{TP+\varepsilon}{TP+FP+\varepsilon}.
\end{equation*}

\section{Worked example}

\code{attention_unet} (published in this catalogue) is a complete method in four
files, about 380 lines total:

\begin{longtable}{@{}p{3.2cm}p{11.2cm}@{}}
\toprule
\textbf{File} & \textbf{Content} \\
\midrule
\endhead
\file{manifest.json} & api 1, \code{train}/\code{predict}/\code{export}, GPL-3 \\
\file{panel.json} & 41 nodes: data, model, training, prediction, export, log --- no Qt code \\
\file{model.py} & U-Net with max-pooling and additive attention gates \\
\file{\_\_init\_\_.py} & handlers, a \code{QObject} bridge, trainer and exporter wiring \\
\bottomrule
\end{longtable}

It is also the ablation baseline for the WRANet comparison, which makes it the
right starting point: copy it, rename the folder and manifest id, and replace
\file{model.py}.

\section{Packaging and publishing}

\subsection{Build the archive}

\begin{lstlisting}[language=bash]
python -m amis.market package ../my_method --catalogue-entry \
    --changelog "What changed in this release"
\end{lstlisting}

It writes \file{my\_method-0.1.0.zip} (single top-level folder, \code{\_\_pycache\_\_}
excluded), prints the SHA-256 and size, and emits the exact \file{index.json}
entry with the real digest.

\subsection{Check it the way the installer will}

\begin{lstlisting}[language=bash]
python -m amis.market verify my_method-0.1.0.zip --id my_method --version 0.1.0
python -m amis.market validate ../my_method --load
\end{lstlisting}

\code{verify} checks the digest, the archive (no absolute paths, no \code{..}, no
symlinks) and the manifest inside. \code{validate} compiles every file, parses
every JSON, and with \code{--load} imports the addon and builds its panel.

\subsection{Publish}

Publishing is a pull request against the catalogue repository:

\begin{enumerate}[leftmargin=1.4em,itemsep=2pt]
  \item upload the ZIP as a release asset (tag
        \code{<id>-v<version>});
  \item append the version to \code{addons[].versions} in \file{index.json};
  \item run \code{python .github/validate\_catalog.py};
  \item open the PR. CI validates, the maintainer merges.
\end{enumerate}

There is no upload button in the application on purpose: every package that
users can install should be reviewed first.

\section{Installing a published addon}

\begin{longtable}{@{}p{4.2cm}p{10.2cm}@{}}
\toprule
\textbf{Channel} & \textbf{How} \\
\midrule
\endhead
Application & \ui{Addon marketplace} \(\to\) \ui{Refresh} \(\to\) select \(\to\) \ui{Install} \\
Terminal & \code{python -m amis.market install <id>} \\
Folder & copy the folder (or a ZIP you trust) into the addons directory \\
\bottomrule
\end{longtable}

Every channel ends the same way: the addon folder appears in the user addons
directory and the registry reloads it.

\section{Security and licensing}

\begin{itemize}[leftmargin=1.4em,itemsep=2pt]
  \item Addons are \textbf{Python code executed in-process}. Installation is a
        trust decision; the marketplace shows version, size and checksum before
        anything is downloaded and refuses a digest mismatch.
  \item Never publish clinical images, patient-derived figures, tokens or
        credentials. Weights trained on restricted data need written permission
        before they appear in a public catalogue.
  \item State your licence. AMIS is GPL-3.0-or-later; an addon installed by the
        user may carry another licence, but anything shipped inside the
        application bundle must be GPL-3 compatible.
  \item Versions are immutable: publish a new version rather than replacing a
        released archive. Mark a bad release \code{"yanked": true} with an
        explanation instead of deleting it.
\end{itemize}

\section{Troubleshooting}

\begin{longtable}{@{}p{4.6cm}p{9.8cm}@{}}
\toprule
\textbf{Symptom} & \textbf{Cause and fix} \\
\midrule
\endhead
Entry shows \textbf{!} with a traceback & import error; read the last line of the
traceback, fix, \ui{Reload addons} \\
\emph{targets API v2} & manifest \code{api} is wrong for this release; set 1 \\
Buttons do nothing & handler name mismatch; use \code{on_<name>_clicked} or set
\code{on_click} \\
\emph{Panel build failed} & the interface raised; validate first, check for an
unknown widget type \\
Widget is grey / missing & a container without \code{children}, or a duplicate
\code{name} \\
Metrics unavailable & no ground truth: the mask must share the image's file stem,
next to it or in \code{masks/} \\
\emph{checksum mismatch} on install & the archive changed after publication;
re-upload as a new version \\
\bottomrule
\end{longtable}

\section{Checklist}

\begin{itemize}[leftmargin=1.4em,itemsep=1pt]
  \item[\(\square\)] folder name equals \code{manifest.id}, \code{api} is 1
  \item[\(\square\)] \code{ADDON} exported; every handler in the panel exists
  \item[\(\square\)] every \code{self.ui.get/set} name exists in the panel
  \item[\(\square\)] widget types are from the table above
  \item[\(\square\)] \code{validate --load} passes; \code{verify} passes
  \item[\(\square\)] licence declared; no data, no secrets
  \item[\(\square\)] numbers formatted with a decimal point
\end{itemize}

\end{document}
